\section{Introduction}
\label{sec:intro}

\subsection{What this monograph follows}

There is a single line of mathematics that runs from the definition of a
Markov decision process to a program that teaches itself to play a board game.
This monograph follows that line in one direction, and then follows the same
line into code.

The mathematical path has six steps:

\begin{enumerate}
  \item the \emph{objective} of a parameterized policy, and the likelihood-ratio
        identity that makes its gradient estimable from samples alone
        (Section~\ref{sec:reinforce});
  \item the \emph{policy gradient theorem}, which removes the derivative of the
        state distribution from the gradient (Section~\ref{sec:pgt});
  \item the \emph{actor--critic} decomposition, and the instability that
        bootstrapping introduces (Section~\ref{sec:actorcritic});
  \item the \emph{deterministic policy gradient}, the sibling branch of the same
        idea, and a clarification of its place in the lineage
        (Section~\ref{sec:dpg});
  \item \emph{approximate policy iteration} and its reformulation as Expert
        Iteration, which is the exact frame around the AlphaZero loop
        (Sections~\ref{sec:apxpi} and~\ref{sec:exit});
  \item the \emph{bandit} results that give tree search its exploration rule and
        its only real guarantees (Section~\ref{sec:bandits}).
\end{enumerate}

The algorithmic path has three steps: the sequence of three DeepMind papers that
deleted one crutch at a time (Section~\ref{sec:three-papers}), the algorithm as a
whole including its search and its loss
(Sections~\ref{sec:algorithm}--\ref{sec:loss}), and the honest inventory of which
parts are proven and which are not (Section~\ref{sec:proven}).

The engineering path has two dominant problems, and this is where most of the
implementation effort actually goes:

\begin{itemize}
  \item \textbf{The shape of the search tree in Rust.} A game tree is a
        recursively owned structure, which is the structure the Rust ownership
        model is least willing to express directly. Section~\ref{sec:tree}
        develops the flat-arena representation, states its invariants, and shows
        the traversal pattern that keeps the borrow checker satisfied without
        runtime checks.
  \item \textbf{The parallel organisation of self-play.} Self-play needs
        thousands of position evaluations per second from one accelerator, while
        the tree search that produces them is irregular, branchy, and cheap. The
        answer is to build many independent trees rather than one shared tree, and
        to make the accelerator a batching service at the single point of
        synchronization. Section~\ref{sec:parallel} develops this.
\end{itemize}

\subsection{The target application}

Every concrete example in this monograph uses freestyle Gomoku on a
$15\times15$ board, with the Swap2 opening protocol. Gomoku is a useful target
for three reasons.

First, it is tactically sharp and short. A game lasts roughly 30 to 70 moves,
and single moves decide games. This is the opposite regime from Go, where
influence accumulates slowly over hundreds of moves. The consequence for the
capital---compute, network size, number of self-play games---is large and
favourable: published AlphaZero-style Gomoku agents reached human playing
strength with networks two orders of magnitude smaller than the Go networks
\citep{xie2018alphagomoku}.

Second, the mechanics are simple enough to be fully specified in a page, which
keeps the rules engine small and keeps this monograph's examples short.

Third, and least obviously, Gomoku has a solved core: Allis and colleagues
proved in 1993 and 1994 that freestyle Gomoku on $15\times15$ is a forced win for
the first player \citep{allis1996gomoku,allis1994thesis}. This is a gift rather
than a problem. Board positions of which the exact value is known, and forced
win sequences that can be verified mechanically, are perfect test data for a
search implementation, and they are available before any learning happens. The
competitive variant, which the Swap2 opening rule produces, is not solved
(Section~\ref{sec:solved}), so the agent still plays on open ground.

\subsection{The ``alpha-epsilon'' position}

A pure AlphaZero agent knows nothing about the game except its rules. This
monograph describes a small, deliberate departure from that purity: hand-written
tactical detection (a win in one move, a forced block, a double threat) computed
by bit operations, together with a bounded threat-space prover.

The departure matters for three reasons, and only the first is about strength.
Tactical detection gives the learning curve a head start, because a young
network's priors are nearly uniform and a prior that already points at immediate
wins saves thousands of fruitless self-play games. The same detection code
serves as a deterministic stand-in evaluator, which allows the entire search to
be built, tested, and proven correct before a single gradient is computed. And
the bounded prover produces labelled positions, whose labels are exact rules
truths rather than bootstrap estimates, which a small fraction of every training
batch can then use as an anchor.

The word \emph{epsilon} in the label is the honest part: the hand-written
knowledge is a small perturbation to an otherwise learned system, not a
replacement for it. Section~\ref{sec:alphaepsilon} states exactly what the
component may and may not do.

\subsection{Scope}

This monograph gives derivations for the mathematics and reasoning for the
implementation. It does not give complete programs, complete hyperparameter
sweeps, or benchmark results, and it does not describe any particular
repository. Code fragments are illustrative and are kept below roughly fifteen
lines where possible.

Three things are deliberately out of scope. Reinforcement learning from
imperfect information (poker, simultaneous games) is not covered, because nothing
in the AlphaZero line applies to it without a substantial reformulation.
Continuous control is covered only at the point where the deterministic policy
gradient enters the narrative (Section~\ref{sec:dpg}); it is discussed because
it is frequently confused with the AlphaZero line, and the confusion is worth
clearing. Finally, distributed training across several accelerators is not
covered: the parallel design in Section~\ref{sec:parallel} is explicitly the
single-accelerator design, and the reasons for the choice are stated there.

\subsection{How to read this monograph}

Readers who want the algorithm should read Sections~\ref{sec:three-papers}
through~\ref{sec:loss} and can treat the mathematical section as reference
material. Readers who want to build the system should read
Sections~\ref{sec:tree} and~\ref{sec:parallel} first: those two sections contain
the design decisions that are expensive to change later.

Readers who are following the original papers should note one recurring
distinction. The primary sources publish algorithms and architectures in full
detail but are reticent about constants. Where a constant is not published, this
monograph says so at the point of use, and marks the substitute as an
experimental choice. Table~\ref{tab:folklore} collects these cases, and this is
the most practically useful table in the document.

All symbols are collected in Appendix~\ref{app:notation} and are redefined at
first use in the body.

\begin{table}[t]
\centering
\small
\caption{Constants commonly attributed to the primary sources but not published
by them. The ``actual source'' column names where the number really comes from.}
\label{tab:folklore}
\begin{tabularx}{\textwidth}{@{}l l X@{}}
\toprule
Constant & Commonly claimed source & Actual source \\
\midrule
$c_{\mathrm{puct}} = 1.5$ & AlphaGo Zero / AlphaZero & Not published by DeepMind;
first widely used published value is from \citet{tian2019elf} \\
Replay window size & AlphaZero & Not published; the 500{,}000-game window is
AlphaGo Zero's, reported via \citet{tian2019elf} \\
Dirichlet $\alpha$ & AlphaZero & The rule $\alpha \approx 10 / (\text{typical legal moves})$
is published; the specific per-game values follow from it \\
Simulations per move & --- & 1600 (AlphaGo Zero), 800 (AlphaZero). Do not mix
them \\
Gating threshold 55\% over 400 games & AlphaGo Zero & Removed entirely in
AlphaZero \\
\bottomrule
\end{tabularx}
\end{table}
